\section*{第七章 总结与展望}
\addcontentsline{toc}{section}{第七章 总结与展望}

\subsection{全文总结}

本文通过 Monte Carlo 数值模拟方法，对五个典型随机过程模型中 OST 的表现进行了系统研究。各模型的核心结论见表~\ref{tab:summary}。

\begin{table}[ht]
  \centering
  \caption{五个模型的 OST 角色对照}
  \label{tab:summary}
  \begin{tabular}{llll}
    \toprule
    模型 & 鞅构造 & OST 角色 & 核心结论 \\
    \midrule
    Moran 模型 & $X_n$ 本身是鞅 & 直接成立 &
    $P(\text{固定}) = x_0/N$ \\
    保险破产 & $e^{-RU_t}$ 指数鞅 & 构造后成立 &
    $\psi(u) \leqslant e^{-Ru}$ \\
    赌徒破产 & $S_n$ 是对称鞅 & 双边成立，单边失效 &
    一致可积性缺失 \\
    秘书问题 & 不是鞅 & 不适用 & Snell 包络 → 37\% 法则 \\
    美式期权 & 折现 Snell 是上鞅 & 结构性保证 &
    LSM 算法定价 \\
    \bottomrule
  \end{tabular}
\end{table}

从表~\ref{tab:summary} 中可以看出 OST 理论谱系的完整面貌：
\begin{enumerate}[noitemsep]
  \item \textbf{OST 成立}（Moran 模型）：过程天然是鞅且停时满足有界条件，OST 直接给出精确等式。
  \item \textbf{OST 通过构造成立}（保险破产）：过程本身不是鞅，但通过指数变换可构造辅助鞅，OST 给出上界而非等式（因停时是单边吸收）。
  \item \textbf{OST 失效}（赌徒破产单边停时）：三个充分条件全部不满足，一致可积性缺失是失效的本质原因。
  \item \textbf{OST 不适用}（秘书问题）：核心过程不是鞅，引入 Snell 包络作为最优停时的统一框架。
  \item \textbf{OST 的结构性角色}（美式期权）：不直接用于导出闭式解，但保证了 LSM 算法的数学合法性。
\end{enumerate}

\subsection{不足与展望}

本文的数值实验基于相对简单的模型设置（指数索赔分布、几何布朗运动等），实际应用中模型可能更为复杂。此外，对于单边停时 OST 失效的理论分析还可以进一步深化——特别是截断停时偏误的收敛速率（$\sim N^{-1/2}$）的理论推导。在数值方法方面，可以探索更高效的方差缩减技术（如对偶变量法、控制变量法）以在相同计算成本下获得更精确的估计。
